\section{Introduction}\label{sec:intro}

\parstart{C}{oughing} is among the few sounds the body produces that carries diagnostic information and
can be recorded without instrumentation, a clinician or a visit. That combination is what
has made automatic cough analysis attractive: a microphone is already in every pocket, and
a cough is loud, brief and distinctive. Almost every application proposed for it, from
monitoring the progression of a chronic respiratory condition to triaging an infectious
one, rests on a step that is rarely the subject of the paper that uses it. Before anything
can be inferred from a cough, the cough has to be found. Detection is the foundation, and
the quality of that foundation sets a ceiling on everything built above it.

The reported accuracy of that foundational step is, on its face, settled. Work on cough
and respiratory audio routinely reports accuracies in the high nineties, on architectures
ranging from purpose built convolutional networks to transfer learning from image
classification~\cite{imran2020, laguarta2020, pahar2021}. A reader encountering these
figures would reasonably conclude that cough detection is solved and that attention
belongs further downstream.

A sequence of independent audits has since established that this conclusion does not
follow from the evidence offered for it. Reviews of machine learning applied to
respiratory audio during the pandemic found that the overwhelming majority of published
models carried a high risk of bias and that almost none were validated on data from a
source other than the one they were developed on~\cite{roberts2021}. More specifically,
models that appeared to detect disease were shown to be reading confounds instead: the
recording device, the collection site, the prompt a participant was reading, or simply the
identity of the speaker, any of which can separate the classes when the classes were not
collected under matched conditions~\cite{han2022, coppock2024}. Where an external test set
was eventually applied, performance fell to near chance~\cite{xia2021, zhou2021}. The
pattern across these audits is consistent, and it is not that the models were badly
trained. It is that the evaluations measured something other than what they were taken to
measure.

This leaves a specific and answerable question, and it is the one this work addresses. The
audits identified that single corpus evaluation inflates performance, but they did not
isolate which part of the evaluation is responsible, because in the studies they examined
several candidate causes vary together. Three in particular are routinely conflated. The
first is the provenance of the label: a recording may be called a cough because a
participant was told to cough, because a person listened and judged that a cough was
present, or because another classifier assigned it a score above a threshold. These are
treated throughout the literature as interchangeable sources of ground truth, and they are
not equivalent claims. The second is the composition of the negative class. Separating
coughing from a sustained vowel is a different problem from separating it from breathing,
sneezing, gasping and laughter, and a corpus whose negatives happen to be easy will
reward a detector that has learned very little. The third is run to run variation.
Differences well under one accuracy point are reported as findings in this area without
any estimate of the spread across random seeds, which is the quantity that determines
whether such a difference can be claimed at all.

The present work holds the architecture, the preprocessing and the training protocol fixed
and varies these three factors deliberately. Three corpora are used, chosen so that each
represents one label regime: prompt derived labels from Coswara~\cite{coswara2020}, human
annotated labels from the exhaustively verified evaluation split of
FSD50K~\cite{fsd50k2022}, and classifier derived labels from the \texttt{cough\_detected}
score shipped with COUGHVID~\cite{coughvid2021}. The class prior is matched across all
three by construction, so a change in performance between regimes cannot be attributed to
a change in class balance. Partitions over the prompt derived corpus are built over
participants rather than recordings, so no speaker appears on both sides of a split.
Every condition is run under five random seeds, and every comparison is tested on paired
per recording decisions rather than on summary accuracies.

The specific contributions are the following.

\begin{itemize}
\item A controlled comparison of the three label regimes under one fixed pipeline, in
      which label provenance is the variable and class prior, preprocessing, architecture
      and training protocol are held constant. This isolates a factor that the existing
      audits identified but did not separate from the others varying alongside it.

\item A quantification of how far single corpus accuracy overstates detection. Under an
      identical input condition, ResNet18 reaches $98.6$ per cent balanced accuracy on
      held out speakers from the prompted corpus and $70.2$ per cent on the human
      annotated split, and the fall is asymmetric in a way that a single accuracy figure
      conceals: specificity declines from $99.5$ to $84.3$ per cent while sensitivity
      collapses from $97.7$ to $56.0$. A detector that appears to be near perfect misses
      close to half of all coughs once the negatives are confusable and the labels come
      from a person. The loss is shown to lie in discrimination rather than in
      calibration, which closes off the cheapest remedy: the area under the curve falls
      from $0.998$ to $0.749$, and moving the decision threshold to the best position for
      the external set recovers only $1.7$ points.

\item An account of what the detectors confuse, obtained from the annotated classes of
      the external negatives. The false positives are concentrated rather than diffuse,
      and they fall on the other forced expirations: clips labelled as sneezes are
      assigned to the cough class in $67.8$ per cent of decisions, a higher rate than the
      same models achieve on clips that are genuinely coughs. The negative class of the
      prompted corpus contains no forced expiration, so the distinction was never one the
      models were required to learn.

\item An assessment of wavelet shrinkage as a preprocessing step under a protocol able to
      detect an effect of the size being claimed for it. Measured against the spread across
      seeds, the gain does not survive: the largest improvement observed anywhere is
      $2.73$ accuracy points, with a spread of $3.54$ points across the seeds it
      is averaged over.

\item A reproducibility finding about the classifier derived corpus. The recording
      selection used by the work this study extends cannot be reproduced from the corpus as
      published, for an arithmetic reason that is demonstrated rather than asserted, and
      the expert annotated subset cannot be used to estimate the corpus label error rate
      because it is not a random sample of the corpus.

\end{itemize}

The computational conditions are stated in enough detail to be acted upon rather than
merely reproduced, and Section~\ref{sec:setup-train} records a measured account of one
failure mode that produces no error and no warning while costing a factor of $6.4$ in
training throughput, together with the measurement that distinguishes its cause from the
plausible alternative.

Nothing in what follows should be read as a claim that cough detection is unachievable, or
that the architectures examined here are inadequate to it. The networks fit the prompted
corpus almost perfectly, and they do so from a speaker disjoint partition, which rules out
the simplest explanation for a high score. The claim is narrower and, for anyone intending
to build on such a result, more consequential: the number obtained that way does not
transfer, the direction in which it fails is predictable, and the experiment needed to
discover this costs far less than the experiment that assumes otherwise.

The remainder of the paper is organised as follows.
Section~\ref{sec:related} reviews classifiers built on cough audio, what has since been
established about how they were evaluated, and the evidence offered for denoising as a
preprocessing step. Section~\ref{sec:design} sets out the pipeline, taking the wavelet
shrinkage, the spectrogram and the networks in the order the signal meets them.
Section~\ref{sec:setup} describes the corpora and their label regimes, the partitioning,
the training protocol and the statistical procedure. Section~\ref{sec:results} reports the
grid, the external evaluation and the comparison across regimes.
Section~\ref{sec:discussion} draws out what follows for how detectors of this kind should
be evaluated and what the corpora would need to contain.
Section~\ref{sec:conclusion} states what the results support and what they do not.
